\documentclass[10pt]{article}

\usepackage[utf8]{inputenc}

\usepackage[T1]{fontenc}

\usepackage{amsmath}

\usepackage{amsfonts}

\usepackage{amssymb}

\usepackage[version=4]{mhchem}

\usepackage{stmaryrd}

\usepackage{bbold}



\begin{document}

стандартного трехмерного куба. Рассматриваются также комплексы, состоящие из бесконечного, но локально конечного семейства многогранников. Понятие П. к. обобщает понятие геометрического симплициального комплекса. Тело $|P| \Pi$. к. $P$ представляет собой объединение всех входящих в него многогранников и является полиэдром. Число многогранников в $P$, как правило, меньше числа симплексов в триангуляции. П. к. $P_{1}$ наз. подразделением комілекса $P$, если их тела совпадают и каждый многогранник из $P_{1}$ содержится в некром многограннике пІз $P .3$ в е 3 д н о е по о д а 3 д ел е н и е комплекса $P$ с центром в точке $a \in|P|$ получается с помощью разбиения замкнутых многогранников, содержащих $a$, на конусы с верщиной $a$ над теми их гранями, к-рые не содөржат $a$. Любой П. к. $P$ имеет подразделение $K$, являющееся геометрическим симплициальным комплексом. Такое подразделение можно получить без добавления новых вершин. Достаточно, напр., последовательно произвести звездные подразделения $P$ с центрами во всех вершинах $P$.



Лum.: [1] А л е к с а н д р о в П. С., Комбинаторная топология, М.- Л., $1947 . \quad$ С. В. Матвеев.



ПОЛИЭДРАЛЬНЫЙ ДИКЛ - полиэдральнал цепь, граница к-рой равна нулю.



C. В. Матвеeв.



ПОЛНАЯ АНАЛИТИЧЕСКАЯ ФУНКЦИЯ - совокупность всех элементов аналитич. функции, получающихся при всевозможных аналитических продолжения $x$ исходной аналитич. функции $f=f(z)$ комплексного переменного $z$, заданной первоначально в нек-рой области $D$ расширенной комплексной плоскости $\overline{\mathbb{C}}$.



Пара $(D, f)$, состоящая из области $D \subset \bar{C}$ и заданной в $D$ однозначной аналитической, или голоморфной, функции $f$, наз. ә л е ме н т о м а на л и т ич еской функции, аналитическия элем е н т о м, или, короче, просто әле м е н то м. Всегда возможно, в частности, при задании аналитич. функции пользоваться в е й е р ш т р а с с о в ы м, или р е г у л я р ны м, ә л е м е н т о м $\left(U(a, R), f_{a}\right)$, состоящим при $a \neq \infty$ из степенного ряда



\[

f_{a}=f_{a}(z)=\sum_{k=0}^{\infty} c_{k}(z-a)^{k}

\]



и круга сходимости $U(a, R)=\{z \in \overline{\mathbb{C}}:|z-a|<R\}$ этого ряда с центром $a$ и радіусом сходимости $R>0$. В случае $a=\infty$ вейерштрассов элемент $\left(U(\infty, R), f_{\infty}\right)$ состоит из ряда



\[

f_{\infty}=f_{\infty}(z)=\sum_{k=0}^{\infty} c_{k^{2}}-k

\]



и области сходимости этого ряда $U(\infty, R)=\{z \in \overline{\mathbb{C}}$ : $|z|>R\}, \quad R \geqslant 0$.



Пусть $E_{f}-$ множество всех тех точек $\zeta \in \overline{\mathbb{C}}$, в к-рые исходный әлемент $\left(U(a, R), f_{a}\right)$ аналитически продолжается хотя бы по одному пути, связывающему в $\overline{\mathbb{C}}$ точки $a$ и $\zeta$. Следует иметь в виду возможность такой ситуации, когда в точку $\zeta \in E_{f}$ аналитич. продолжение возможно вдоль нек-рого класса путей $L_{1}$ и невозможно вдоль другого класса путей $L_{2}$ (см. Особая точ$\kappa a$ аналитич. функции). Множество $E_{f}$ есть область плоскости $\overline{\mathbb{C}}$. П о л н й ан алит и ч е с ко й фу кци е й (в с мы с ле В е й е р ші т а с с а) $f_{W}$, порожденной элементом $\left(U(a, R), f_{a}\right)$, называется совокупность всех вейерштрассовых элементов $\left(U(\zeta, R), f_{\zeta}\right), \zeta \in E_{f}$, получаемых при этом аналитич. продолжении вдоль всевозможных путей $L \subset \overline{\mathbb{C}}$. Область $E_{f}$ наз. (в е й е р ш т р а с с о в о й) о б л ас т в ю с у щ е с т в о в а ни я П. а. ф. $f_{W}$. Применяя элементы общего вида $(D, f)$, вместо вейерштрассовых на самом деле получают ту же самую П. а. ф. $f_{W}$. Элементы $(D, f)$ П. а. Ф. $f_{W}$ часто наз. ветвями аналитической функиии $f_{W}$. Любой элемент $(D, f)$ П. а. ф. $f_{W}$, будучи взят ва исходный при аналитич. продолжении. приводит к той же самой П. а. ф. $f_{W}$ Каждый элемент $\left(U(\zeta, R), f_{\zeta}\right)$ П. а. ф. $f_{W}$ может быть получен из любого другого ее элемента $\left(U(a, R), f_{a}\right)$ посредством аналитич. продолжения вдоль нек-рого пути, связывающего в $\overline{\mathbb{C}}$ точки $a$ и $\zeta$.



Может оказаться, тто исходный әлемент $(D, f)$ не допускает аналитич. продолжения ни в одну точку $\zeta \notin D$. В этом случае $D=E_{f}$ является е с т е с т в е нн о й обла с т в 10 с ущ е с т во в а н и я, или облас т ю г голо м о рфности, функции $f$, а ее граница $\Gamma=\partial D-$ e с т е с т в е н о й $г \mathrm{p}$ а нице й функции $f$. Напр., для вейерштрассова әлемента



\[

\left(U(0,1), f_{0}(z)=\sum_{k=0}^{\infty} z^{k !}\right)

\]



естественной границей является окружность $\Gamma=\{z \in$ $\in \overline{\mathbb{C}}:|z|=1\}$ его круга сходимости $U(0,1)$, т. к. атот әлемент нельзя продолжить аналитически ни в одну точку $\zeta$ такую, что $|\zeta| \geqslant 1$. Какова бы ни была область $D \subset \overline{\mathbb{C}}$, можно построить аналитич. функцию $f_{D}(z)$, для к-рой $D$ есть естественная область существования $f_{D}(z)$, a ee граница $\Gamma=\partial D-$ естественная граница $f_{D}(z)$ (это следует, напр., из Миттаг-Лебблера теоремьи).



П. а. Ф. $f_{W}$ в своей области существования $E_{f}$, вообще говоря, не является функцией точки в обычном смысле этого слова. Часто встречающаяся в теории аналитич. функций ситуация такова, что П. а.ф. $f_{W}$ есть многозначная функция: для каждой точки $\zeta \in E_{f}$ существует, вообще говоря, бесконечное множество элементов $\left(U(\zeta, R), f_{\zeta}\right)$ с центром в этой точке. Однако это множество не более чем счетное (т е о р ем а П у а нка р е - В оль т е р р а). В целом П. а. ф. $f_{W}$ можно рассматривать как однозначную аналитич. функцию только на соответствующей римановой поверхности, являющейся многолистной накрывающей поверхностью над $\overline{\mathbb{C}}$. Напр., П. а. ф. $f(z)=$ $=\operatorname{Ln} z=\ln |z|+i \operatorname{Arg} z$ многозначна в своей области существования $E_{f}=\{z \in \overline{\mathbb{C}}: 0<|z|<\infty\}$; в каждой точке $\zeta \in E_{f}$ она принимает счетное множество значений



$f_{\zeta}(\zeta ; s)=\ln |\zeta|+i \arg \zeta+2 \pi s i, s=0, \pm 1, \ldots$,



и каждой точке $\zeta \in E_{f}$ соответствует счетное множество элементов



\[

\begin{gathered}

\left(U(\zeta,|\zeta|), f_{\zeta}(z ; s)\right)= \\

=f_{\zeta}(\zeta ; s)+\sum_{k=1}^{\infty} \frac{(-1)^{k-1}}{k \zeta^{k}}(z-\zeta)^{k}

\end{gathered}

\]



с центром $\zeta$. Обычно используется однозначная ветвь әтой П. а.ф. - главное значение логарифма $\ln z=$ $=\ln |z|+i \arg z$, являющееся голоморфной функцией в области $D=\{z \in \overline{\mathbb{C}}: 0<|z|<\infty,-\pi<\arg z<\pi\}$, непрерывно продолжаемой на множество $\{z \in \overline{\mathbb{C}}: 0<$ $<|z|<\infty,-\pi<\arg \xi \leqslant \pi\}$.



При обращении (см. Обращение ряда) вейерштрассовых элементов (1), (2) возникают элементы более общей природы, определяемые соответственно р яд а м и П юи и ё:



\[

f_{a}=\sum_{k=\mu}^{\infty} c_{k}(z-a)^{k / v}, \quad f_{\infty}=\sum_{k=\mu}^{\infty} c_{k^{z-k / v}},

\]



где $\mu$ - целое число, $v$ - натуральное число, и кругами сходимости этих рядов $U(a, R), U(\infty, R)$. В частности, при $\mu \geqslant 0, \quad v=1$ ряды (3) совпадают с рядами (1), (2), определяющими регулярные элементы; в отличие от них определяемые рядами (3) элементы при $\mu<0$ или $v>1$ наз. о с о б ы м и. При $\boldsymbol{v}=1$ и $v>1$ ряды (3) определяют соответственно н е р а з в е т в-





\end{document}